\chapter{Đứa Phá Lớp Không Chỉ Là Phiền}
\markboth{Đứa Phá Lớp Không Chỉ Là Phiền}{The Disruptive Student Is Not Only a Problem}

\section{Gây Ồn Để Khỏi Bị Xem Thường}
\begin{truyen}{Gây Ồn Để Khỏi Bị Xem Thường}{Being Loud to Avoid Feeling Small}
\chuhoa{C}{ó em hễ vào lớp là muốn phải có người nhìn tới mình. Không bằng chuyện tốt thì bằng chuyện ồn.}
\textit{(Some students need to be noticed the moment they enter the room. If not through something good, then through noise.)}

Nhiều khi đằng sau cái ồn đó là một nỗi sợ rất cũ: sợ mình không đáng kể.
\textit{(Behind that noise is often an old fear: the fear of being insignificant.)}

Khi không biết làm mình có giá trị bằng cách nào khác, các em dễ chọn cách gây chú ý.
\textit{(When a child does not know how else to feel valuable, attention-seeking becomes an easy path.)}

Cái phiền trên mặt có khi chỉ là một nỗi thiếu bên trong.
\textit{(What looks annoying on the surface may simply be an inner lack.)}

\end{truyen}

\begin{baihoc}
Không phải lúc nào tiếng ồn cũng đến từ hư. Có khi nó đến từ cảm giác mình quá nhỏ.
\textit{(Noise does not always come from badness. Sometimes it comes from feeling too small.)}
\end{baihoc}

\begin{ghinhoanh}
\textbf{Short English to remember:}

Noise can come from feeling too small.

\end{ghinhoanh}

\begin{tuhocnhanh}
\textbf{Gợi ý để dạy và tự nghĩ / Teaching and reflection prompts}

1. Tiếng ồn này có thể đang nói hộ điều gì? \textit{What might this noise be speaking for?}

2. Em học sinh này đang thiếu cảm giác gì? \textit{What feeling might the student be missing?}

3. Mình có thể cho em một chỗ đứng khác không? \textit{Can we offer the student a different kind of place?}
\end{tuhocnhanh}
\ngancach

\section{Chọc Bạn Để Che Điểm Yếu}
\begin{truyen}{Chọc Bạn Để Che Điểm Yếu}{Teasing Others to Hide Weakness}
\chuhoa{C}{ó em rất nhanh miệng. Hễ bạn nào sai là em chêm vào một câu làm cả lớp cười.}
\textit{(Some students are sharp-tongued. The moment another student slips, they jump in with a joke that makes the class laugh.)}

Nhìn lâu mới thấy em thường làm vậy nhất vào những lúc mình không theo kịp bài.
\textit{(Watch long enough and you notice it happens most when the student is not keeping up.)}

Chọc người khác giúp em giấu đi cảm giác yếu của chính mình.
\textit{(Mocking others helps the student hide personal weakness.)}

Đôi khi một đứa trẻ làm tổn thương người khác chỉ vì nó đang cố cứu cái tôi của mình.
\textit{(Sometimes a child hurts others simply to rescue a shaky ego.)}

\end{truyen}

\begin{baihoc}
Không phải hành vi xấu nào cũng bắt đầu từ ác ý. Nhiều khi nó bắt đầu từ xấu hổ.
\textit{(Not every bad behavior begins with cruelty. Many begin with shame.)}
\end{baihoc}

\begin{ghinhoanh}
\textbf{Short English to remember:}

Many hurtful behaviors begin with shame.

\end{ghinhoanh}

\begin{tuhocnhanh}
\textbf{Gợi ý để dạy và tự nghĩ / Teaching and reflection prompts}

1. Hành vi này đang che điều gì? \textit{What is this behavior covering?}

2. Xấu hổ dẫn tới làm đau người khác bằng cách nào? \textit{How can shame lead to hurting others?}

3. Người dạy nên xử lý cả hành vi lẫn phần nào bên trong? \textit{Should the teacher address both the behavior and the inner cause?}
\end{tuhocnhanh}
\ngancach

\section{Phá Nhưng Lại Rất Thương Thầy}
\begin{truyen}{Phá Nhưng Lại Rất Thương Thầy}{Disruptive Yet Deeply Attached}
\chuhoa{C}{ó em trong giờ thì nghịch, nói chen, làm thầy cô mệt. Nhưng đến khi thầy bệnh hoặc nghỉ dạy, em lại là đứa hỏi đầu tiên.}
\textit{(Some students disrupt class and exhaust the teacher, yet they are the first to ask when the teacher is sick or absent.)}

Cái cách các em thương nhiều khi vụng và rất khó chịu.
\textit{(The way they care is often clumsy and difficult.)}

Không phải đứa trẻ nào thương người khác cũng biết cách cư xử cho dễ thương.
\textit{(Not every caring child knows how to behave in a pleasant way.)}

Nhìn học trò, nhiều lúc phải phân biệt giữa bản chất và biểu hiện.
\textit{(With students, one often has to distinguish between the heart and the behavior.)}

\end{truyen}

\begin{baihoc}
Có những đứa trẻ tốt theo cách rất vụng.
\textit{(Some children are good in very clumsy ways.)}
\end{baihoc}

\begin{ghinhoanh}
\textbf{Short English to remember:}

Some children are good in very clumsy ways.

\end{ghinhoanh}

\begin{tuhocnhanh}
\textbf{Gợi ý để dạy và tự nghĩ / Teaching and reflection prompts}

1. Điều gì trong em này đáng giữ lại? \textit{What in this student is worth preserving?}

2. Vì sao cần tách bản chất khỏi biểu hiện? \textit{Why should heart and behavior be distinguished?}

3. Người dạy sửa phần nào trước? \textit{What should the teacher address first?}
\end{tuhocnhanh}
\ngancach

\section{Bị Dán Nhãn Quá Sớm}
\begin{truyen}{Bị Dán Nhãn Quá Sớm}{Labeled Too Early}
\chuhoa{N}{hiều em chỉ cần một vài tháng đầu nghịch nhiều là bị cả lớp, cả giáo viên, thậm chí cả chính mình tin rằng em chỉ là “đứa phá.”}
\textit{(Some students misbehave for a few months and soon the whole class, the teachers, and even the student begin to believe that “troublemaker” is the whole identity.)}

Từ đó về sau, mỗi lần em cố tốt lên cũng khó được nhìn lại từ đầu.
\textit{(After that, every attempt to improve struggles to receive a fresh look.)}

Nhãn dán sai không chỉ làm người lớn nhìn lệch. Nó còn làm trẻ con sống theo cái lệch đó.
\textit{(A wrong label does not only distort adult perception. It pushes the child to live inside that distortion.)}

Có những cơ hội sửa mình chết từ một cái tên gọi quá sớm.
\textit{(Some chances to change die inside a label given too early.)}

\end{truyen}

\begin{baihoc}
Đừng để một giai đoạn trở thành cả căn cước của một đứa trẻ.
\textit{(Do not let one phase become a child’s entire identity.)}
\end{baihoc}

\begin{ghinhoanh}
\textbf{Short English to remember:}

Do not let one phase become a child’s entire identity.

\end{ghinhoanh}

\begin{tuhocnhanh}
\textbf{Gợi ý để dạy và tự nghĩ / Teaching and reflection prompts}

1. Cái nhãn này đang làm hại điều gì? \textit{What is this label damaging?}

2. Vì sao trẻ dễ sống theo nhãn bị gắn? \textit{Why do children often live according to the label given?}

3. Người lớn có thể mở lại cơ hội bằng cách nào? \textit{How can adults reopen a chance?}
\end{tuhocnhanh}
\ngancach

\section{Tới Giờ Kiểm Tra Thì Im Thin Thít}
\begin{truyen}{Tới Giờ Kiểm Tra Thì Im Thin Thít}{Completely Quiet When the Test Starts}
\chuhoa{C}{ó em bình thường phá nhất, nhưng tới giờ kiểm tra thì ngồi im như một người khác.}
\textit{(Some students are the most disruptive in normal lessons, yet become completely silent during tests.)}

Điều đó cho thấy em không hề không biết sợ, không hề không biết giữ.
\textit{(That shows the student does know how to restrain and focus.)}

Nhiều em phá không phải vì không còn gì để giữ. Các em chỉ chưa biết giữ mình ở đâu cho đúng.
\textit{(Many disruptive students are not empty of self-control. They simply do not yet know where to place it well.)}

Cái đó đáng phạt, nhưng cũng đáng dạy.
\textit{(That may deserve discipline, but it also deserves teaching.)}

\end{truyen}

\begin{baihoc}
Nếu một học sinh biết giữ ở lúc này, nghĩa là em vẫn còn cái để xây ở lúc khác.
\textit{(If a student can hold themselves together in one setting, there is still something solid to build on elsewhere.)}
\end{baihoc}

\begin{ghinhoanh}
\textbf{Short English to remember:}

If a student can hold together once, there is still something to build on.

\end{ghinhoanh}

\begin{tuhocnhanh}
\textbf{Gợi ý để dạy và tự nghĩ / Teaching and reflection prompts}

1. Điều gì được lộ ra trong giờ kiểm tra? \textit{What becomes visible during the test?}

2. Vì sao điều đó đáng để hy vọng? \textit{Why is that worth hoping on?}

3. Mình có đang nhìn ra phần còn xây được không? \textit{Are we seeing the part that can still be built?}
\end{tuhocnhanh}
\ngancach

\section{Gặp Riêng Thì Rất Lễ Phép}
\begin{truyen}{Gặp Riêng Thì Rất Lễ Phép}{Very Respectful in a One-on-One Moment}
\chuhoa{C}{ó những em trong đám bạn thì làm dữ, nói cứng, nhìn rất ngang.}
\textit{(Some students act hard, speak roughly, and look defiant in front of friends.)}

Nhưng gọi ra riêng thì lại “dạ, thưa” rất rõ.
\textit{(But when spoken to privately, they answer with clear respect.)}

Lúc đó mới thấy nhiều hành vi xấu của học trò mang màu của đám đông nhiều hơn màu của bản chất.
\textit{(That reveals how much some bad behavior belongs more to the crowd than to the child’s core nature.)}

Muốn sửa một số em, đôi khi phải kéo các em ra khỏi khán giả trước.
\textit{(To correct some students, we sometimes must remove the audience first.)}

\end{truyen}

\begin{baihoc}
Có những đứa trẻ hư hơn khi có người xem.
\textit{(Some children behave worse when they have an audience.)}
\end{baihoc}

\begin{ghinhoanh}
\textbf{Short English to remember:}

Some children behave worse when they have an audience.

\end{ghinhoanh}

\begin{tuhocnhanh}
\textbf{Gợi ý để dạy và tự nghĩ / Teaching and reflection prompts}

1. Đám đông đang làm em này thành người như thế nào? \textit{How is the crowd shaping this student?}

2. Vì sao nói riêng có hiệu quả hơn? \textit{Why can private conversation work better?}

3. Người dạy cần chọn bối cảnh sửa ra sao? \textit{How should the teacher choose the setting for correction?}
\end{tuhocnhanh}
\ngancach

\section{Ở Nhà Chẳng Ai Nghe Em Nói}
\begin{truyen}{Ở Nhà Chẳng Ai Nghe Em Nói}{No One Listens at Home}
\chuhoa{C}{ó em làm ồn trong lớp vì đó là nơi hiếm hoi em nói mà có người đáp lại.}
\textit{(Some students become noisy in class because it is one of the few places where speaking gets any response.)}

Ở nhà, có khi em bị ngắt, bị mắng, hoặc chẳng ai có thời gian nghe trọn một câu.
\textit{(At home, the child may be interrupted, scolded, or simply unheard.)}

Nên vào lớp, em nói như người sợ mất lượt.
\textit{(So at school the student speaks like someone afraid of losing their turn.)}

Không phải hành vi nào cũng sinh ra trong chính lớp học.
\textit{(Not every behavior is born inside the classroom itself.)}

\end{truyen}

\begin{baihoc}
Muốn hiểu tiếng ồn của học sinh, đôi khi phải nghe cả phần im lặng ở nhà các em.
\textit{(To understand student noise, we sometimes have to hear the silence waiting for them at home.)}
\end{baihoc}

\begin{ghinhoanh}
\textbf{Short English to remember:}

To understand student noise, we may need to hear the silence at home.

\end{ghinhoanh}

\begin{tuhocnhanh}
\textbf{Gợi ý để dạy và tự nghĩ / Teaching and reflection prompts}

1. Tiếng ồn này có thể bắt đầu từ đâu? \textit{Where might this noise begin?}

2. Vì sao nhà và lớp nối với nhau? \textit{Why are home and classroom connected?}

3. Người dạy cần hiểu sâu đến mức nào? \textit{How deeply should the teacher try to understand?}
\end{tuhocnhanh}
\ngancach

\section{Cho Em Một Việc Thật, Em Bớt Phá}
\begin{truyen}{Cho Em Một Việc Thật, Em Bớt Phá}{Give the Student a Real Job and the Noise Often Drops}
\chuhoa{T}{ôi từng giao một em hay phá việc lau bảng, phát phiếu, nhắc bạn xếp bàn ghế.}
\textit{(I once gave a disruptive student the jobs of cleaning the board, handing out papers, and helping remind classmates to arrange the room.)}

Không phải ngày nào em cũng tốt lên ngay.
\textit{(The change did not happen perfectly or immediately.)}

Nhưng rõ ràng em phá ít hơn khi cảm thấy mình có ích.
\textit{(But the student definitely disrupted less when feeling useful.)}

Có những đứa trẻ không cần thêm mắng. Chúng cần một chỗ để thấy mình có giá trị.
\textit{(Some children do not need more scolding. They need a place where they can feel useful.)}

\end{truyen}

\begin{baihoc}
Trách nhiệm nhỏ đôi khi kéo một đứa trẻ ra khỏi vai “đứa phiền.”
\textit{(A small responsibility can sometimes pull a child out of the role of “the annoying one.”)}
\end{baihoc}

\begin{ghinhoanh}
\textbf{Short English to remember:}

A small responsibility can pull a child out of the role of “the annoying one.”

\end{ghinhoanh}

\begin{tuhocnhanh}
\textbf{Gợi ý để dạy và tự nghĩ / Teaching and reflection prompts}

1. Vì sao việc nhỏ lại có tác dụng? \textit{Why can a small responsibility help?}

2. Em này đang cần cảm giác gì? \textit{What feeling does this student need?}

3. Mình có thể giao việc gì cụ thể? \textit{What concrete task can we assign?}
\end{tuhocnhanh}
